% BHU PYQ -- Auto-generated & error-corrected
\documentclass[11pt,a4paper]{article}

\usepackage[a4paper,top=2.5cm,bottom=2.5cm,left=2.8cm,right=2.8cm,headheight=14pt]{geometry}
\usepackage{amsmath,amssymb}
\usepackage[T1]{fontenc}
\usepackage[utf8]{inputenc}
\usepackage{lmodern}
\usepackage{microtype}
\usepackage{fancyhdr}
\usepackage{enumitem}
\usepackage{hyperref}
\usepackage{lastpage}
\usepackage{parskip}

\hypersetup{colorlinks=true,urlcolor=black,linkcolor=black,
  pdftitle={BSC-07A: Descriptive Statistics -- BHU PYQ},pdfauthor={Banaras Hindu University}}

\pagestyle{fancy}\fancyhf{}
\renewcommand{\headrulewidth}{0.4pt}\renewcommand{\footrulewidth}{0.4pt}
\fancyhead[L]{\small   \textit{Banaras Hindu University}}
\fancyhead[R]{\small   \textit{BSC-07A: Descriptive Statistics}}
\fancyfoot[C]{\small Page  \thepage\ of \pageref{LastPage}}


\newlist{parts}{enumerate}{2}
\setlist[parts,1]{label=(\alph*),leftmargin=*,itemsep=5pt,topsep=3pt}
\setlist[parts,2]{label=(\roman*),leftmargin=*,itemsep=3pt,topsep=2pt}

\newcommand{\pts}[1]{\hfill   \textit{[\#1]}}


\begin{document}


\begin{center}
  {\large\bfseries BANARAS HINDU UNIVERSITY}\\[2pt]
  {
\normalsize B.Sc. (Hons.) Semester II Examination, 2023-24}\\[6pt]
  \rule{0.8\linewidth}{0.5pt}\\[6pt]
  {\large\bfseries Paper No. BSC-07A: Descriptive Statistics}\\[4pt]
  {
\normalsize Time: 3 Hours\hfill Maximum Marks: 70}
\end{center}
\rule{\linewidth}{0.4pt}
\smallskip



\noindent   \textit{Instructions: Answer   \textbf{five} questions including Question~1 which is compulsory. Symbols have their usual meanings.}
\bigskip

MRL ING scsascsacinresttaacrevexcsctiices
BSC-O7A : Statistics - I: Descriptive Statistics -
\medskip


\noindent   \textbf{Question 1.}


\begin{parts}
  \item Define primary and secondary data sets. Also mention its important sources.
  \item What are the different types of the measurement scale?Describe each with suitable example.
\end{parts}
\medskip\rule{\linewidth}{0.2pt}

\medskip


\noindent   \textbf{Question 2.}


\begin{parts}
  \item What do you mean by ungrouped frequency distribution? In an opinion survey, the household size of different families living in Gokul Dham Colony of Varanasi district is recorded as:
  
  2, 2, 3, 3, 4, 4, 5, 5, 6, 2, 2, 3, 4, 5, 6, 3, 4, 4, 4, 4, 4, 4, 5, 2, 2, 3, 4, 4
  
  Construct the frequency distribution for the above data set. Also construct the cumulative and relative frequency tables. \pts{7}
  
  \item What are the different types of classification available to represent the raw data set? Explain with suitable example. \pts{7}
\end{parts}
\medskip\rule{\linewidth}{0.2pt}

\medskip


\noindent   \textbf{Question 3.}


\begin{parts}
  \item Describe the meaning and utility of histogram and Ogive.
  \item Define quartiles, deciles and percentiles?Compute different quartiles for the Income in Rs. (in 000) | 5-10 10-15 20-25 25-30 No. of Workers 10 \pts{16}
\end{parts}
\medskip\rule{\linewidth}{0.2pt}

\medskip


\noindent   \textbf{Question 4.} What are the various mathematical and positional averages?Describe each with their

following data set.
merits and demerits.
\medskip


\noindent   \textbf{Question 5.} What are the different measures of dispersion?Describe each in detail. Also discuss the

effect of change of origin and scale on variance.
\medskip


\noindent   \textbf{Question 6.} Define correlation with its various properties. What are the different assumptions are

required for defining correlation?The Marks obtained by the group of 10 students in
the subject Statistics (X) and Mathematics (Y) out of 20 are given below. Obtain the
Pearson's coefficient of correlation between X and Y.
Statistics (X) 6 9 es \pts{15}
Mathematics (Y) | 9 3 Seah S
(

. Define regression and regression coefficient. Write all the properties of regression
coefficient. Obtain the regression lines for the given data set.
(X, Y): (5,10), (7,5), (4,8), (8,4), (10,5)
. Write sort notes on any four of the followings.
i. Bar diagram and Pie chart
ii. | Weighted mean
iii. Coefficient of dispersions
iv. Skewness and Kurtosis
v. Raw and central moments
-o 

\vfill
\rule{\linewidth}{0.4pt}

\begin{center}\small
  \href{https://bhu-exam.vercel.app/}{Click here to visit our website}\\[4pt]\href{https://me-aryan.vercel.app/}{Click here to visit our contributor}\\[4pt]\href{https://quashan-me.vercel.app/}{Click here to visit our contributor}
\end{center}

\end{document}